\subsection{Arc graphs in brane labelled windowed surfaces}

A {\it windowed surface} $F=F^s_g(\delta_1,\ldots ,\delta _r)$ is a
smooth oriented surface of genus $g\geq 0$ with $s\geq 0$ punctures and
$r\geq 1$ boundary components together with the specif\/ication of a
non-empty f\/inite subset $\delta _i$ of each boundary component,
for $i=1,\ldots ,r$, and we let $\delta =\delta
_1\cup\cdots\cup\delta _r$ denote the set of all distinguished points in the boundary
$\partial F$ of $F$ and let  $\sigma$ denote the set of all punctures.
The set of components of
$\partial F-\delta$ is called the set $W$ of {\it windows}.

Furthermore one needs to specify a brane labelling $\beta$.
For this we f\/irst f\/ix a set $\calB$ of basic brane labels
and denote by $\calP(\calB)$  power set.  Notice that $\varnothing\in\calP(\calB) $,
this will encode the closed sector. The elements of $\calP(\calB)$ of cardinality
bigger than one should be thought of as intersecting branes. These
give room for extra data, but it is possible to set all the contributions
for these  ``higher intersection'' branes
to zero in a given model.


A {\it brane-labeling} on a windowed surface $F$ is a function
\[
\beta : \ \delta\coprod\sigma\to{\mathcal P}({\mathcal B}),
\]
where $\sqcup$ denotes the disjoint union, so that if $\beta
(p)=\varnothing$ for some $p\in\delta$, then $p$ is the unique point
of $\delta$ in its component of~$\partial F$.  A brane-labeling
may take the value $\varnothing$ at a puncture.



A window $w\in W$ on a windowed surface $F$ brane-labeled by
$\beta$ is called {\it closed} if the endpoints of $w$ coincide at
the point $p\in\delta$ and $\beta (p)=\varnothing$; otherwise, the
window $w$ is called {\it open}.
Each window def\/ines a pair of brane labels $\beta(w)$ which is
the pair $(S,T)$ of the brane labels of the beginning and end of the window (these
may coincide).

\subsubsection{Arc families}

We def\/ine the sets
\[
\delta
(\beta )=\{ p\in\delta :\beta (p)\neq\varnothing\},\qquad \sigma
(\beta )=\{ p\in\sigma :\beta (p)\neq\varnothing\} .
\]


Def\/ine a $\beta$-arc $a$ in $F$ to be an arc properly embedded in $F$ with its endpoints in $W$ so that~$a$ is not
homotopic, f\/ixing its endpoints, to $\partial F-\delta (\beta )$.  For example, given a distinguished point
$p\in\partial
F$, consider the arc lying in a small neighborhood that simply connects one side of $p$ to another in $F$; $a$ is
a $\beta$-arc if and only if
$\beta (p)\neq\varnothing$. We will call such an arc a {\it small arc} around $p$.

Two $\beta$-arcs are {\it parallel} if they are homotopic rel~$\delta$, and a $\beta$-arc family is the homotopy
class
rel~$\delta$ of a collection of $\beta$-arcs, no two of which are parallel.  Notice that we take homotopies rel~$\delta$
rather than
rel $\delta (\beta )$.





Given a positively weighted arc family in $F$, let us furthermore say that a window $w\in W$ is {\it active}
if there
is an arc in the family with an endpoint in $w$, and otherwise the window is {\it inactive}.



\subsubsection{The mapping class group and arc graphs}

The {\it $($pure$)$ mapping class group} $MC(F)$ of $F$ is the group of
orientation-preserving homeomorphisms of $F$ pointwise f\/ixing
$\delta\cup\sigma$ modulo homotopies pointwise f\/ixing $\delta\cup\sigma$.


$MC(F)$ acts naturally on the set of $\beta$-arc families.
An {\it arc graph} is an equivalence graph under this action.

\subsection{Arc spaces, moduli spaces and the open/closed Sullivan PROP}
\subsubsection{The Arc spaces}



A {\it  weighting} on an arc family is the assignment of
a positive  real number to each of its components.
A weighting naturally passes to the arc graph.
A weighting is called {\it discrete} if it takes values in ${\mathbb N}$.



Let ${\rm Arc}'(F,\beta )$ denote the geometric realization of the partially
ordered set of all $\beta$-arc families in $F$.  ${\rm Arc}'(F,\beta )$ is
described as the set of all projective positively weighted $\beta$-arc
families in $F$ with the natural topology. (See for instance \cite{KLP} or \cite{P1} for
further detail.)

$MC(F)$ again
acts naturally on ${\rm Arc}'(F,\beta )$.  The {\it arc
complex} is def\/ined to be the  quotient under this action
\[
{\rm Arc}(F,\beta )={\rm Arc}'(F,\beta )/MC(F).
\]

We shall also consider the corresponding deprojectivized versions:
$\widetilde{\rm Arc}{}'(F,\beta )\approx {\rm Arc}'(F,\beta )\times{\mathbb R}_{>0}$ is the space of all positively weighted arc
families in $F$ with the natural topology, and
\[
\widetilde{\rm Arc}(F,\beta )=\widetilde{\rm Arc}{}'(F,\beta )/MC(F)\approx
{\rm Arc}(F,\beta )\times{\mathbb R}_{>0}.
\]



For any windowed surface $F$, def\/ine
\[
\widetilde{\rm Arc}(F)=\bigsqcup \widetilde{\rm Arc}(F,\beta),
\]
where the disjoint union is over all brane-labellings on $F$.
\[
\widetilde{\rm Arc}(n,m)=\bigsqcup \left\{ \alpha\in \widetilde{\rm Arc}(F):
\begin{array}{l}\alpha ~{\rm has}~n~{\rm closed~and}~ m~{\rm
open~active}\\{\rm
windows~and~no~inactive~windows}\end{array}\right\},
\] where the
disjoint union is over all orientation-preserving
homeomorphism classes of windowed surfaces.




